\documentclass[11pt,a4paper]{article}

\title{\textbf{Business Requirements Document (BRD)}\\[0.5cm] \textbf{AeroGrid}\\[0.3cm] Low-Altitude Airspace Intelligence Warehouse}
\author{AeroGrid Development Team}
\date{\today}

\begin{document}

\maketitle

\section{Document Purpose}
This Business Requirements Document (BRD) details the stakeholder needs, operational business processes, analytical requirements, and operational constraints for the AeroGrid project.

\section{Business Aims \& Analytical Focus}
The warehouse serves both day-to-day monitoring and long-term planning by addressing the following core needs:
\begin{itemize}
    \item Delivering a unified operational awareness view for low-altitude airspace traffic.
    \item Quantifying airspace utilization across altitude layers and geographic grids.
    \item Detecting, locating, and analyzing close encounters (near-midair conflicts).
    \item Evaluating the operational impact of regulatory restrictions and adverse weather.
    \item Informing strategic expansion and corridor design for emerging drone operations.
\end{itemize}

\section{Business Processes}
The warehouse is structured around five core business processes representing actual operations in low-altitude airspace:

\begin{enumerate}
    \item \textbf{Continuous Vehicle Tracking:} Captures ongoing position reports (position, height, speed, heading) broadcast by active aircraft and drones, forming the foundation for real-time awareness and density calculations.
    \item \textbf{Complete Flight Journey Recording:} Records the entire path of each journey (takeoff to landing or monitored area entry/exit) to analyze route choices, duration, altitude profiles, and traffic flow.
    \item \textbf{Restriction Activation and Coverage:} Records standing and temporary flight restrictions (TFRs around airports, stadiums, security sites), including effective time windows and geographic boundary polygons.
    \item \textbf{Close Encounter Detection:} Identifies and logs instances where vehicles breach minimum safe separation distances in space and time, turning raw telemetry into measurable safety events.
    \item \textbf{Airspace Occupancy Snapshot:} Calculates periodic utilization metrics (vehicle counts, occupancy duration per grid cell, busiest altitude bands) to support capacity planning.
\end{enumerate}

\section{Scope \& Constraints}
\begin{itemize}
    \item \textbf{Data Scope:} Publicly available surveillance feeds (OpenSky), FAA UAS datasets, TartanAviation trajectories, and weather services.
    \item \textbf{Users:} Traffic management service providers, regulatory authorities, commercial drone operators, and municipal city planners.
    \item \textbf{Technical Constraints:} Designed using standard dimensional data warehousing principles (Star/Snowflake schemas) optimized for complex analytical queries.
\end{itemize}

\end{document}
